\textbf{18.} На основании теоремы Рейнгольда ($\mathsf{UPATH} \in \mathbf{L}$) докажите принадлежность к $\mathbf{L}$ множества связных графов, множества циклических графов и множества двудольных графов.

\textbf{Решение.} Напомним, что сей язык определяется так:
\[
\mathsf{UPATH} = \{(G, s, t)\colon \text{в неориентированном графе $G$ есть путь из $s$ в $t$}\}.
\]

Связность проверяется совсем легко: давайте создадим по два счётчика, которые будут символизировать собой две вершины графа $s$ и $t$, и для каждой такой пары запустим алгоритм на логарифмической памяти, проверяющий наличие пути. Если в графе всего $n$ вершин, то размер входа уже хотя бы $n$, а значит хранить два счётчика размером $\log{n}$ мы сможем.

В проверке, есть ли цикл в графе, напрашивается следующий алгоритм: перебрать все рёбра (то есть пары смежных вершин, счётчики по ним будут логарифмическими), в таком цикле удалить рассматриваемое ребро $e = (u, v)$ из графа и проверить, есть ли путь между вершинами $u$ и $v$ в новом графе. Если цикл был, то для соседних вершин из цикла путь остался, и наоборот. Единственная техническая трудность заключается в том, как именно удалить ребро --- входную ленту, на которой дан граф, менять нельзя. Давайте тогда всякий раз, когда алгоритм проверки наличия пути пытается обратиться к матрице смежности графа, проверять, наличие какого именно ребра он хочет проверить. Если он обращается к биту, соответствующему паре $(u, v)$, мы будем давать ему 0, иначе --- то, что в действительности лежит на входной ленте.

Вспомним, что двудольность графа равносильна отсутствию циклов нечётной длины. Используем стандартную технику: <<раздвоим>> вершины. Если в старом графе была вершина $u$, то в новом графе ей соответствуют две вершины $u_1$ и $u_2$. Также для ребра $(u, v)$ в старом графе мы добавим рёбра $(u_1, v_2)$ и $(u_2, v_1)$ в новом. Нужно это для того, чтобы при проходе по ребру чётность индексов менялась. Теперь если мы посмотрим на алгоритм проверки наличия цикла из предыдущего пункта задачи, то заметим, что цикл нечётный, если и только если путь, соединяющий вершины $u$ и $v$ в новом графе без ребра $(u, v)$, -- чётной длины. Но чётность пути с помощью раздвоенного графа проверяется легко: должен быть путь между $u_1$ и $v_1$ (то есть вершинами одной чётности). Таким образом, нам надо также пройтись в цикле по всем рёбрам графа, <<удалить>> ребро, раздвоить вершины и проверить, есть ли путь между вершинами одной чётности. Конечно же, возникают те же сложности, что и ранее: мы не можем явно построить раздвоенный граф --- памяти нет. Поэтому придётся считать его сходу: если алгоритм проверки наличия пути спрашивает, есть ли ребро между вершинами графа, мы по их номерам восстанавливаем чётность и номера вершин в исходном графе, по которым находим ответ на запрос.